\section{Evaluation protocol}
\label{sec:protocol}

\subsection{Pre-registration workflow}
Every confirmatory comparison reported in this paper follows the same sequence. (1)~A protocol is written and committed to the repository before the individuals under test are trained: the systems, the training command, the evaluation worlds (including their random seeds), the hypotheses with effect-size thresholds, the statistics and the known limitations. (2)~The individuals are trained and none is selected, excluded or retrained. (3)~The SHA-256 hashes of all evolved genomes are written and committed. (4)~The evaluation script, which refuses to run if any hash differs from the committed one, is run \emph{once}. (5)~Results are recorded as they are, including those that contradict the hypotheses. Git commits supply the time stamps. Pre-registrations are numbered 02--11 in the repository, in the order in which they were written. Development-phase experiments, which informed the design of later protocols, are labelled as such and are not used as confirmatory evidence.

Two procedural deviations are disclosed. In pre-registration 04 the stepwise rule did not anticipate tied $p$ values (Section~\ref{sec:results}), and from pre-registration 06 the Holm step-down uses $p$ values only, with the effect-size condition applied to each hypothesis separately. In pre-registration 09 the hash commit was made after, not at, the end of training (the code and the genomes are unchanged, verified by the hashes of the source files recorded in each training manifest), and evaluation was split into eight shards run in parallel, checked against a single-process run on 36 worlds.

\subsection{Unit of analysis and statistics}
The unit of analysis is the \emph{training seed}: the same evolutionary noise sequence is used for every condition, so a comparison is a paired difference between conditions that share a seed. For each seed we average fitness over the evaluation worlds (two deception levels with 32 worlds each, over one or two evaluation pools) and take the difference between the two conditions. With $n$ seeds, the test of a positive mean difference is the exact one-sided sign-flip test: the $p$ value is the fraction of the $2^n$ sign assignments whose mean is at least the observed mean, so with $n=8$ the smallest attainable value is $1/256\approx 0.0039$. Where the claim is non-inferiority, the differences are shifted by the margin first. Hypotheses within a pre-registration are corrected with the Holm step-down at a total level of 0.05. A hypothesis is \emph{supported} only if it is rejected by the Holm procedure \emph{and} its mean difference meets its own pre-registered threshold. Confidence statements about magnitude use percentile bootstrap intervals over the training seeds (20,000 resamples); with three to eight seeds these intervals are rough, and we report them only as descriptions.

Pre-registrations 02 and 03 predate this design: they resample worlds within seeds (02) or require significance in each of three training seeds (03). We keep their formal outcomes and report them as stated.

\subsection{Tiered reading of results}
A single pass/fail verdict conflates direction and size. We therefore also report, after the fact, a tiered reading of each comparison (Appendix, Table~\ref{tab:tiers}): \emph{direction confirmed} when all seeds have the same sign; \emph{magnitude confirmed} when the lower bound of the seed-level interval reaches the pre-registered threshold, \emph{magnitude unclear} when the interval contains it, and \emph{magnitude below threshold} when the upper bound lies below it; \emph{non-inferiority confirmed} when the lower bound exceeds the margin. This reading does not change any formal outcome. Its purpose is to prevent two errors: reading ``not supported'' as ``no effect'' when a small effect is certain in sign, and reading ``supported'' as ``large'' when the interval is wide.

\subsection{Evaluation pools and reuse}
Evaluation uses held-out dialect pools whose descriptions never appear in training. The held-out field-notes pool was evaluated in several pre-registrations (only summary statistics were inspected; it was never used to tune a design or to select a genome), so it is no longer an untouched sample. Where a new comparison needs an untouched pool, we use the factual-dialect pool, which was first used for a prior-quality evaluation and then for individuals in pre-registration 06. Hypothesis tests are not corrected across pre-registrations; in total 228 independent evolutionary runs contribute to the confirmatory comparisons (12 in pre-registrations 02--03, 24 in 04, 48 in each of 06 and 07, and 24 in each of 08--11).
